%% 致谢

\begin{acknowledgements}
一直憧憬着毕业时的轻松与自由，真正坐下来提笔写致谢的这一天，才发觉一切比想象中沉重得多。当然，对我而言，未来还有博士要读，前路未竟，多少也是一重因素。不过更多的是因为，回忆从来不是轻松的。它像压缩饼干，将四年的时光紧紧凝结其中：哭过，笑过，有成长，有收获。真到了分别这一天，不舍是自然的，也是应当的。

在写这段的时候，想挑首歌听，脑海里第一个浮现的，是五月天的《如烟》。"有没有那么一首诗篇，找不到句点，青春永远定格在我们的岁月"。青春之所以宝贵，恰恰在于时间从不为任何人驻足，也更不可能回退。此刻站在人生的阶段交接点上，青春如烟四个字，有了格外真切的重量。

这首歌第一次听，是大一的时候在电影艺术鉴赏课上，姚经纬老师课上分享给我们的。姚老师是我大学里最喜欢的老师之一。您幽默风趣，与学生打成一片，讲课却又鞭辟入里、见解深刻，是我电影艺术的启蒙人，让我真正发现了电影作为第七艺术的那一面。您与我们年岁相近，亦师亦友，不仅在学业上给予我启发，更在生活上帮我解决了许多大学之前从未有人告诉我该如何面对的问题。大学的重心固然是学习，但人不能只活在学习里。在那些被课业填满的日子里遇见您，我深感幸运。谨在此，致以最诚挚的感谢。

说到艺术，我也想郑重地感谢大艺团合唱队这个大家庭，同时，也借此正式道一声别。在合唱队待了将近两年。刚加入时，疫情与种种因素的影响下，全队不过十几个人，可以说是在合唱队最艰难的时候入的队。初印象就是有一个悉心教我发声、识谱的廖队，有一个有时候很凶有时候有趣的刘丹老师，还有一堆志同道合热爱音乐的伙伴。我们一起排练，一起登台，度过了许多充实的时光，也收获良多。最看得见的收获，是排练场上结交的朋友，是那许多有趣的灵魂，是跟着老师和时常造访我们学校的音乐大师们学到的点点滴滴。而更深处无形的收获，是让我培养了我认为最重要的，对音乐的那份热爱。

我也想感谢张凯龙教授，您的计算机组成原理课堂非常生动有趣。您深知内驱力才是学习最好的老师，总是积极与学生互动，调动我们的主动性。印象最深刻的莫过于您分享的汪国真的一首散文诗，其中的那句"世界有不绝的风景，我有不老的心情"，至今仍是我的个性签名。您跟我们说要对一切保持好奇，抱持学习的心态，永远怀揣探索未知世界的欲望。我想，做科研也是如此。作为一名准博士生，这种探索欲将会是我坚持下去的最核心的源动力，并以此为基础去发现、去创新、去突破。

我想感谢实验室的各位老师、师兄师姐与同门。首先，诚挚感谢我的导师郭斌教授。您作为我科研路上的引路人，在我刚刚踏入实验室、对科研一无所知的时候，是您带我入门，不仅在选题、研究方法和论文写作上给予了我细致而耐心的指导，更重要的是让我真正理解了什么是科研。过去我对科研的想象或许还停留在"提出方法、做出结果、写成论文"的层面，但在您的指导下，我逐渐明白，科研更像是一场长期的自我训练：训练自己如何发现问题，如何面对失败，如何在无数次推翻和重来之后仍然保持清醒与耐心。您严谨的治学态度、开放的学术视野以及对学生成长的关心，都会成为我今后继续走在科研道路上的重要指引。能够继续跟随您攻读博士，是我本科阶段做出的最重要的决定之一，我深感荣幸，也深知责任重大。

同时，感谢实验室的各位师兄师姐，尤其是丁亚三师兄、李林炜师兄、程如意师姐、林政师兄。你们在我科研起步阶段给予了我莫大的帮助，无论是指导实验方法，还是在组会上为我的汇报提出宝贵意见，抑或是你们分享的一路走来的经验等，都让我感受到了一个团队该有的温度。也要感谢与我同届的同门们，我们一起学习，一起成长，讨论问题、合作项目、交流经验，也一起吐槽、分享趣事、嘻嘻哈哈。你们是我科研路上最好的伙伴！

此外，感谢大学四年里遇到的所有同学与朋友。是你们，给了我多元而美好、值得铭记一辈子的大学记忆。也许我们会走向各自不同的未来，也许彼此都只能有一段阶段性的陪伴，但我们在一起度过的时光永远会流淌在记忆里，是我最珍贵的宝藏。

最后，我想感谢我的家人。你们总是倾尽所有给予我，永远无条件地站在我身后。无论我想做什么，选择哪条路，就算你们未必完全看懂我的想法，也总是选择包容与理解。也许成长意味着离家越来越远，意味着很多事情都要独自面对，但我始终知道，身后有一个地方，会在我疲惫时接住我。我爱你们。

很喜欢卢广仲的《几分之几》，它有一个浪漫而奇妙的英文名，叫做《You Complete Me》。我想，这大概也是相遇的意义。阶段性的陪伴，其实是人生中最常见的状态；对于"永远到底有多远"这个问题，没有人敢轻易回答。人与人相遇、碰撞，摩擦、交织、离别，就算是最后离开了对方的生命，还是会在一定程度，一定方面，永永远远留下影子，大概就是被你改变的那部分。我，代替你永远陪在了我的身边。所以所有在生命中相遇的人，谢谢你成为我的几分之几。
\end{acknowledgements}
